%% IEICE Technical Report (ieicej.cls) - dense 8-page version without external figs
%% NOTE: This class requires pLaTeX/uplatex workflow (DVI -> PDF).
%% For pdflatex, use alternative packages

\documentclass[technicalreport]{ieicej}
\ifx\pdftexversion\undefined
  \usepackage[dvipdfmx]{graphicx}
\else
  \usepackage{graphicx}
\fi

\usepackage[fleqn]{amsmath}
\usepackage{newtxtext}
\usepackage[varg]{newtxmath}
\usepackage{url}

\urlstyle{rm}

% ===== Optional macros =====
\newcommand{\LLM}{LLM}
\newcommand{\GCN}{GCN}

% ===== Metadata =====
\tecrepyear{2026}
\titlepagebaselinestretch{0.82}

\jtitle{機能的フローネットワークと気象情報を用いた \GCN-\LLM\ Agent による鉄道旅客需要変動の説明生成}
\etitle{Explaining Railway Passenger Demand Fluctuations with a \GCN-\LLM\ Agent Using Functional Flow Networks and Weather Data}

\authorlist{%
 \authorentry{成瀬 宥哉}{Yuya NARUSE}{Oita}
 \authorentry{大知 正直}{Masanao OCHI}{Oita}
}

\affiliate[Oita]{大分大学 理工学部 共創理工学科 知能情報システムコース}
{Oita University, Faculty of Science and Technology, Intelligent Information Systems Course}

\begin{document}

\begin{jabstract}
鉄道旅客需要予測では，数値的な需要変動だけでなく，その増減要因を運用担当者が理解可能な形で提示することが重要である。本研究では，駅間共変動に基づく機能的フローネットワークと気象情報を統合し，\GCN\ による需要推定と \LLM\ Agent による根拠付き説明生成を接続する枠組みを提案する。評価では，\GCN\ 特徴量により需要推定誤差（MSE）が0.5331から0.3673へ改善した。また，気象情報を根拠として与えることで，降雨・週末の要因検出で高いF1（0.718, 0.696）を得た。以上より，予測精度と説明の妥当性を同時に扱う枠組みの有用性を示す。
\end{jabstract}

\begin{jkeyword}
LLM，LLM Agent，グラフ畳み込みネットワーク，交通需要予測，説明生成
\end{jkeyword}

\begin{eabstract}
Railway passenger demand forecasting requires not only numerical forecasts but also explanations of why demand increases or decreases. This paper proposes a framework that connects demand estimation with \GCN\ and grounded explanation generation with an \LLM\ Agent by integrating a functional flow network and weather data. The functional flow network is constructed from co-movement patterns among stations, and weather conditions and neighboring-station changes are provided to the \LLM\ Agent as evidence candidates. Experiments show that adding \GCN\ features reduced the MSE from 0.5331 to 0.3673. The explanation task also achieved high F1 scores for rain and weekend factors, 0.718 and 0.696, respectively. These results indicate that the proposed framework can support both prediction accuracy and evidence-based explanation generation.
\end{eabstract}

\begin{ekeyword}
LLM, LLM Agent, GCN, passenger demand, explanation generation
\end{ekeyword}

\maketitle

\section{はじめに}
都市鉄道では，ダイヤ調整，混雑緩和，要員配置などの運用判断を短期的な旅客需要予測に基づいて行う。交通系ICカードの普及により，駅単位・時間単位の乗降データを用いた需要分析が可能となり，LSTMやTransformerなどの系列モデル，さらに時空間GNNを用いた交通需要予測が発展してきた\cite{7508408,vaswani2017attention,10.5555/3304222.3304273}。一方で，実運用では予測値そのものに加えて，「なぜ需要が増減したのか」を運用担当者が理解できる形で提示することが求められる。

本研究が対象とする鉄道旅客需要の説明生成には，次の三つの課題がある。第一に，単一駅の時系列のみでは，路線網上の移動や乗換を介した駅間波及を十分に扱えない。RNNやLSTMは時間依存の表現に有効であるが\cite{7508408,sherstinsky2020fundamentals,yu2019review}，鉄道需要では観測点である駅同士がネットワークとして結び付いているため，時間依存と駅間依存を同時に扱う必要がある。第二に，降雨，気温，風速，週末性などの外生要因と，駅間の共変動に基づく依存関係を，予測精度向上だけでなく説明の根拠としても利用可能にする必要がある。時空間GNNや適応的グラフ学習は交通予測の高精度化に寄与しているが\cite{li2017diffusion,10.5555/3304222.3304273,9363197,wu2019graph,bai2020adaptive}，得られた構造を運用者が解釈できる説明へ接続する設計は十分ではない。第三に，\LLM\ を用いた自然言語説明では，一般知識に基づく補完が混入し，観測された数値事実や鉄道固有のネットワーク構造と整合しない説明が生成されるリスクがある。交通流予測に \LLM\ を用いる研究は現れ始めているものの\cite{guo2024towards}，どの外部要因とどの駅間情報を根拠として与えるかを明示的に設計することが不可欠である。

そこで本研究では，気象要因と駅間依存を明示的な根拠として扱い，\GCN\ による需要推定と \LLM\ Agent による説明生成を統合した枠組みを提案する。具体的には，ICカード乗降データから駅間の共変動に基づく機能的フローネットワークを構築し，\GCN\ により近傍駅の変動を集約した特徴量を得る。さらに，降雨量や気温などの外生要因と近傍駅変動を候補根拠として \LLM\ Agent に入力し，需要変動の要因を自然言語で出力する。これにより，数値予測と説明生成を分離せず，予測段階で得られた構造情報を説明の根拠として再利用できる。

本研究の貢献は以下の三点である。第一に，駅間共変動に基づく機能的フローネットワークを構成し，物理的隣接だけでは捉えにくい需要連動を \GCN\ 特徴量として需要推定へ組み込む。第二に，気象要因と近傍駅変動を \LLM\ Agent の候補根拠として与え，観測事実に基づく説明生成を行う。第三に，要因ラベルを付与可能な模擬データを用いて，予測性能だけでなく，生成説明の妥当性をHitRate，macro-F1，NeighborAcc，ROUGE-L，BERTScore等により定量評価する。
\section{関連研究}
本節では，前節で述べた三つの課題に対応して，(1) 時系列需要予測とネットワーク構造，(2) 駅間依存・外生要因の統合，(3) 根拠付き説明生成と評価，の観点から関連研究を整理する。

\subsection{時系列需要予測とネットワーク構造}
交通需要予測の基盤として，RNNやLSTMに代表される系列モデルが広く用いられてきた。LSTMは長期依存の学習困難を緩和する構造として提案され\cite{7508408}，その性質や応用可能性も整理されている\cite{sherstinsky2020fundamentals,yu2019review}。また，Transformerは自己注意機構により系列内の長距離依存を並列的に扱う枠組みを確立した\cite{vaswani2017attention}。一方，交通流は時間変化だけでなく空間的な相互作用を伴うため，早くから時空間モデルの重要性が指摘されてきた\cite{kamarianakis2005space}。鉄道需要においても，駅を独立した時系列として扱うだけでは，乗換駅やハブ駅を介した需要変動を十分に捉えられない。

\subsection{駅間依存と外生要因の統合}
駅間依存を扱う手法として，グラフ構造上で特徴を集約する \GCN\ が基礎的枠組みを与えている\cite{kipf2016semi,defferrard2016convolutional}。交通予測では，DCRNNが拡散畳み込みと系列モデルを組み合わせて道路ネットワーク上の方向性を表現し\cite{li2017diffusion}，STGCNは時空間畳み込みにより交通流予測の性能を高めた\cite{10.5555/3304222.3304273}。さらに，AST-GCNは属性情報を統合することで，外部情報を含む交通予測の有効性を示している\cite{9363197}。一方で，固定された隣接行列に依存しない方法として，Graph WaveNetやAGCRNでは適応的な隣接関係をデータから学習する手法が提案されている\cite{wu2019graph,bai2020adaptive}。また，公共交通需要は天候の影響を受けることが報告されており\cite{zhou2017impacts}，気象条件を外生要因として扱うことは重要である。しかし，これらの研究の多くは予測精度の改善を主目的としており，駅間依存や気象要因を説明生成時の根拠として明示的に提示する点には十分踏み込んでいない。

\subsection{根拠付き説明生成と評価}
説明可能性の観点では，注意機構やXAIの利用が検討されている。GATは近傍ノードへの重みを学習できるため，どのノードが予測に寄与したかを示す手がかりとなりうる\cite{velickovic2017graph}。また，公共交通需要予測において深層学習と説明可能AIを組み合わせる研究も行われている\cite{monje2022deep}。近年は，交通流予測に \LLM\ を用いて自然言語説明を付与する研究も報告されており\cite{guo2024towards}，数値予測を運用者に理解しやすい説明へ変換する可能性が示されている。ただし，\LLM\ 単体では観測事実に基づかない補完が混入する可能性があり，説明生成に用いる根拠を明示的に設計する必要がある。さらに，生成説明の評価には表層一致を測るROUGE\cite{lin2004rouge}や，意味的類似度を測るBERTScore\cite{scorato2023bertscore}などが用いられるが，鉄道需要説明では要因ラベルの一致や近傍駅波及の検出といったドメイン固有の評価も必要となる。

以上より，時系列予測，時空間GNN，外生要因統合，説明可能AIおよび \LLM\ による説明生成はそれぞれ発展してきた。しかし，駅間共変動に基づく機能的フローネットワーク，気象要因，\GCN\ による近傍集約特徴を統合し，それらを根拠として \LLM\ Agent による説明生成へ接続し，さらに説明の妥当性を定量評価する枠組みは十分に整備されていない。本研究はこの不足を補うことを目的とする。

\section{提案手法}
本研究は、ICカード乗降データに対して、周期性、外部要因、駅間依存の三つの信号を統合し、需要変動の予測と要因説明を一体化する自動生成フレームワークを提案する。
\subsection{全体アーキテクチャ}
本研究では，予測と説明を一体化した枠組みを提案する。既存のアプローチでは，予測モデルが内部的に駅間依存を学習する一方で，その学習結果を運用者に説明する仕組みが十分整備されていなかった。本提案では，\GCN\ で得られた隠れ特徴量を，\LLM\ Agent の入力となる根拠情報へ直接変換することで，予測の背景にある因果構造を自然言語で明示的に表現することを試みる。

提案フレームワークは，(1) 入力データ整備，(2) 機能的フローネットワーク構築，(3) \GCN\ による需要推定，(4) \LLM\ Agent による根拠付き説明生成，の4段階から成る。
\begin{figure}[t]
  \centering
  \fbox{\parbox{0.90\linewidth}{\centering
    \Large 本研究の全体アーキテクチャ図\\[1em]
    （図1：提案フレームワークの概要）\\
    入力データ整備$\rightarrow$機能的フローネットワーク$\rightarrow$\GCN\ による需要推定$\rightarrow$\LLM\ Agent\ による根拠付き説明生成
  }}
  \caption{提案フレームワークの全体アーキテクチャ}
  \label{fig:overview}
\end{figure}

\subsection{入力データと前処理}
入力として，駅ごとに1時間単位で集計したICカード乗降データを用いる。実験では154時間分のログと1,451駅（首都圏2次ODデータ）を対象とした。需要系列の裾の重さを緩和するため，駅$i$，時刻$t$における乗降者数 $x_{i,t}$ に対して次式の対数変換を行う。
\begin{equation}
y_{i,t}=\log(1+x_{i,t})
\label{eq:log1p}
\end{equation}
その後，各駅ごとにZスコア正規化を施し，駅規模の差による支配を抑える。時間情報は時刻・曜日を $\sin,\cos$ による循環特徴量として与え，23時から0時への不連続を回避する。これにより，モデルは絶対規模ではなく増減パターンの類似性に着目しやすくなる。

\subsection{外部要因信号の構成}
外部要因として，降雨量，気温，風速，週末フラグ等を扱う。これらは予測モデルの特徴量として用いるだけでなく，\LLM\ Agent に対して「当該時刻にどの要因が成立しうるか」を示す根拠情報として入力する。説明生成では，単なる数値列ではなく，要因候補として解釈可能な形式に整形し，例えば「降雨あり」「高温条件成立」「近傍駅で直前急増あり」といった入力へ変換する。

\subsection{機能的フローネットワークの構築}
駅間依存は，単純な地理隣接ではなく，共変動に基づく機能的結合として表現する。駅$\times$時間帯行列から駅ペアの類似度を計算し，各駅について相関上位$k$近傍のみを残すことで疎な隣接行列 $A$ を構成する。学習安定化のため，自己ループ付きの対称正規化行列
\begin{equation}
\hat{A}=D^{-1/2}(A+I)D^{-1/2}
\label{eq:adj}
\end{equation}
を用いる。これにより，路線図上の接続だけでなく，利用パターン上の強い結び付きも捉えられる。

\subsection{\GCN\ による需要推定}
需要推定では，時間特徴のみの線形モデル（Linear only）と，\GCN\ 由来の近傍集約特徴を追加したモデル（Linear + GCN-feats）を比較する。\GCN\ は予測器というより，機能的フローネットワーク上で近傍情報を集約する特徴抽出器として利用する。損失関数には，駅平均利用量 $\mu_i$ に基づく重み付きMSEを用いる。
\begin{equation}
L=\frac{1}{T}\sum_{t\in T}\frac{1}{N}\sum_{i=1}^{N} w_i(\hat{y}_{i,t}-y_{i,t})^2,
\quad
w_i=\frac{1}{\sqrt{1+\mu_i}}
\label{eq:wmse}
\end{equation}
これにより，大規模ハブ駅の誤差に学習が過度に引っ張られることを防ぐ。需要予測の精度向上だけでなく，どの近傍が影響しているかという説明の土台も同時に得られる。ここで使用する近傍集約特徴は，式\eqref{eq:adj}で定義される正規化隣接行列 $\hat{A}$ を用いて定義される。

\subsection{\LLM\ Agent による根拠付き説明生成}
\LLM\ Agent には，(1) 対象駅・対象時刻の予測増減量，(2) 降雨量・気温・風速・週末フラグ等の気象状態，(3) 近傍駅の直前変動の要約統計を根拠候補として与える。Agent は許可ラベル集合 $\mathcal{L} = \{$降雨，高温，強風，週末・祝日，近傍駅波及$\}$ を参照し，各根拠候補と予測増減との因果関係を評価した上で最も整合するラベルを選択し，その理由を1〜2文で記述した説明レポートを出力する。自由生成に任せず許可ラベル集合で選択肢を限定することで，\LLM\ の内部知識だけに依存した根拠不明の補完を抑制し，観測事実との整合性を高めている。

\subsection{提案手法の特徴}
本提案手法の最大の特徴は，予測モデルと説明生成器を密に統合することにある。既存の説明可能AI (XAI) では，予測モデルの内部状態を後から可視化する手法が主流であり，その説明は往往いば抽象的・一般論的になりがちである。一方，本研究では，\GCN\ によって抽出された近傍特徴を直接 \LLM\ Agent に供給することで，「駅Aで急増したため，隣接駅Bにも波及した」といった具現的かつ因果的な説明を可能にする。この設計により，運用担当者が直ちに検証可能な根拠を伴った説明が得られ，予測値の信頼性向上と運用判断の精度向上に寄与する。

\subsubsection{データフローの概要}
提案フレームワークにおけるデータフローは，図\ref{fig:overview}のように，単純なデータ処理から構成される。ICカード乗降ログはまず，対数変換と正規化を経て時系列特徴として得られる。これらの特徴量は，\GCN\ を通じて近傍駅の情報を集約し，隠れ特徴ベクトルへと変換される。この隠れ特徴ベクトルは，気象特徴量と合わせて \LLM\ Agent に根拠として提示され，最終的に自然言語で説明が生成される。このように，数値信号 → グラフ特徴 → 説明文章という一貫したデータフローを確立することで，各段階の変換が追跡可能であり，運用担当者が説明の信頼性を検証できる体制が整う。

\section{実験}
本章では，提案手法の有効性を，予測性能と説明性能の両面から検証する。評価は三つの観点から行った。第一に，\LLM\ Agent に外部要因を与えることで説明性能が改善するかを確認した。第二に，機能的フローネットワークと \GCN\ により需要推定精度が向上するかを評価した。第三に，気象情報と駅間波及情報の統合が，説明の妥当性にどのように寄与するかをアブレーションで検証した。

\subsection{説明評価用模擬データ}
正解ラベルが既知の評価環境を構築するため，ベース需要 $B_{s,t}$ に対して外部要因の影響を重畳した模擬データを生成した。要因集合を $\mathcal{K}$，要因$k$の影響度を $w_k$，時刻$t$の要因特徴を $x_{k,t}$ とすると，観測需要を次式で与える。
\begin{equation}
y_{s,t}=B_{s,t}\left(1+\sum_{k\in\mathcal{K}}w_kx_{k,t}\right)+\varepsilon_{s,t}
\label{eq:synth}
\end{equation}
ここで $\varepsilon_{s,t}$ は観測ノイズである。これにより，各要因の寄与を既知とした状態で説明性能を評価できる。また，駅間波及については，式\eqref{eq:synth}の要因集合 $\mathcal{K}$ を構成した上で，基準系列に対する近接シグナルを与え，neighbor ラベルの正解を構成した。

\subsection{予測対象データ}
需要推定では，154時間分の時間別ログと1,451駅（首都圏2次ODデータ）を対象とした。目的変数は駅ごとの乗降者数であり，式\eqref{eq:log1p} により対数変換した後，駅単位でZスコア正規化を施した。さらに，時間特徴として時刻・曜日の循環特徴量を与え，周期性を明示した。これにより，高利用駅と低利用駅が混在する条件でも，増減パターンの学習がしやすくなる。

\subsection{説明生成の比較条件}
説明生成の比較条件は以下の4通りである。
\begin{itemize}
\item BASE：系列履歴のみを用いる
\item +WEATHER：気象要因を追加する
\item +GCN：近接駅波及情報を追加する
\item +WEATHER+GCN：気象要因と駅間波及の両方を追加する
\end{itemize}

また，需要推定では以下の3手法を比較した。
\begin{itemize}
\item Last-hour persistence
\item Ridge (L=6)
\item Ridge + GCN-feats
\end{itemize}

\subsection{評価指標}
需要推定はMSEで評価する。説明生成については，以下の3つの観点から評価する。

まず，要因分解の妥当性として，正解ラベル集合と生成説明から抽出される要因集合の一致に基づき，HitRate，macro-F1，NeighborAccを用いる。HitRate は正解要因を1つ以上含む説明の割合，macro-F1 は要因集合間のF1スコアの平均，NeighborAcc は正解ラベルに近傍波及が含まれる事例における近傍波及ラベルの選択正解率を表す。

さらに，説明文の品質評価として，ROUGE-L~\cite{lin2004rouge} と BERTScore~\cite{scorato2023bertscore} を追加する。ROUGE-Lはn-gramの重複度を評価し，参照説明との表層的類似度を測る。BERTScoreは事前学習済み言語モデルの文脈表現を用いて意味的整合性を評価する。

加えて，時刻$t$における旅客数 $N_t$ を用いた変動度
\begin{equation}
V_t=\frac{N_t-N_{t-7}}{N_{t-7}}
\label{eq:vt}
\end{equation}
を定義する。ここで $N_{t-7}$ は7時間前の旅客数である。生成説明から抽出した各要因のスコアと $V_t$ のスピアマン順位相関をパターン一致度として算出し，変動の強さと説明要因との整合性を評価した。

\subsection{機能的フローネットワークの診断}
機能的フローネットワークの妥当性を確認するため，ノード数，エッジ数，平均次数に加え，エッジ上の平均類似度，ラグ付き相互相関，最大相関が $\pm1$ band に収まる比率を測定した。さらに，次数上位ノードを除去した hub ablation により，主要駅欠損時の頑健性も確認した。

\subsection{実装設定}
需要推定では，log1p変換後のデータを時間順に train (110時間) / test (44時間) へ分割し，将来データで評価した。Ridge回帰ではL=6の自己回帰特徴に加え，駅単位の時刻・曜日の循環特徴量を用いた。Ridge + GCN-feats では，これに式\eqref{eq:adj}で定義される正規化隣接行列 $\hat{A}$ により近傍加重平均したGCN特徴量（2次元，前期間と当期）を追加した。説明生成では，Qwen2.5-7B-Instruct を \LLM\ として用い，許可ラベル集合と候補根拠を入力し，要因ラベルの選択と短い要約生成を行わせた。これにより，自由生成よりも根拠との整合性を保ちやすくした。

\section{結果}
本章では，機能的フローネットワークの構造妥当性，需要推定性能，および \LLM\ Agent による説明生成性能を順に示す。特に，\GCN\ 特徴量の追加が需要推定誤差に与える影響，ならびに気象要因と駅間波及情報の提示が \LLM\ による説明の妥当性にどの程度寄与するかに着目して結果を整理する。

\subsection{機能的フローネットワークの妥当性}
提案フレームワークに基づき構築した機能的フローネットワークは，1451ノード，13243エッジ，平均次数18.25となった。エッジ上の平均類似度は0.868，ラグ付き相互相関の平均は0.869であり，最大相関が $\pm1$ band に収まる比率は0.998であった。これは，結合駅間の変動がほぼ同期的であり，\GCN\ の近傍集約仮定と整合することを示している。

\begin{table}[t]
\caption{機能的フローネットワークの要約}
\label{tab:graph}
\begin{center}
\begin{tabular}{l|r}
\Hline
指標 & 値 \\
\Hline
ノード数 & 1451 \\
エッジ数 & 13243 \\
平均次数 & 18.25 \\
平均類似度 & 0.868 \\
mean lagged xcorr & 0.869 \\
frac($|\mathrm{lag}|\le1$) & 0.998 \\
\Hline
\end{tabular}
\end{center}
\end{table}

さらに，次数上位ノードを除去した hub ablation では，除去後も残余エッジの平均類似度が高い水準を保ち，ネットワーク全体の共変動構造への影響が限定的であることを確認した。これは，鉄道網の需要連動が特定の主要駅のみに依存せず，複数の代替的な共変動経路を持つことを示唆している。

\subsection{需要推定結果}
表\ref{tab:pred}に需要推定結果を示す。Ridge 回帰 (L=6) は Last-hour baseline と同等の予測性能を示すが，GCN-feats を追加することで MSE を 0.5331 から 0.3673（31.1\% 改善）に低減した。これは，時間的特徴に加えて駅間依存を近傍特徴として加えることが有効であることを示す。

\begin{table}[t]
\caption{需要推定結果（MSE，小さいほど良い）}
\label{tab:pred}
\begin{center}
\begin{tabular}{l|r}
\Hline
手法 & MSE \\
\Hline
Last-hour persistence & 0.5075 \\
Ridge (L=6) & 0.5331 \\
Ridge + GCN-feats & \textbf{0.3673} \\
\Hline
\end{tabular}
\end{center}
\end{table}

この結果は，単駅の履歴だけでは捉えにくい短期変動が，近傍駅の同時変動を加えることで補足できることを意味する。特に，Ridge回帰単体ではLast-hour単純持続に及ばないものの，GCN特徴量の追加により大幅な改善が得られた点は，ネットワーク構造の情報が需要推定において本質的な役割を果たすことを示している。

\subsection{説明生成結果}
表\ref{tab:exp3}に \LLM\ Agent による説明生成結果（Qwen2.5-7B-Instruct使用）を示す。気象・近傍情報の両方を与えた条件（+WEATHER+GCN）において，HitRate は 58.1\%，macro-F1 は 0.298 を達成した。特に降雨と週末の要因検出では高いF1スコア（降雨: 0.718, 週末: 0.696）を示しており，外生要因の明示的入力が説明性能に寄与していることが確認された。一方，高温（F1: 0.000）や強風（F1: 0.077），近傍駅波及（NeighborAcc: 0.0\%）については検出が困難であり，これらの要因を \LLM\ が数値データのみから因果関係として認識するにはより明示的な特徴量設計が必要であることを示唆している。

\begin{table}[t]
\caption{\LLM\ Agent による説明生成結果（+WEATHER+GCN条件）}
\label{tab:exp3}
\begin{center}
\begin{tabular}{l|c}
\Hline
指標 & 値 \\
\Hline
HitRate & 0.581 \\
macro-F1 & 0.298 \\
F1(降雨) & 0.718 \\
F1(週末) & 0.696 \\
F1(強風) & 0.077 \\
F1(高温) & 0.000 \\
F1(近傍駅波及) & 0.000 \\
NeighborAcc & 0.000 \\
\Hline
\end{tabular}
\end{center}
\end{table}

\begin{table}[t]
\caption{主要結果の要約}
\label{tab:main}
\begin{center}
\begin{tabular}{l|l}
\Hline
需要推定（MSE） & \begin{tabular}{@{}l@{}}0.5075 (Last-hour)\\
$\rightarrow$ 0.3673 (Ridge+GCN)\end{tabular} \\
説明（+WEATHER+GCN条件） & \begin{tabular}{@{}l@{}}HitRate 0.581\\
macro-F1 0.298\end{tabular} \\
\Hline
\end{tabular}
\end{center}
\end{table}

観測結果を条件ごとに見ると，+WEATHER+GCN条件では，降雨（F1: 0.718）や週末（F1: 0.696）のように数値閾値と需要変動の対応関係が明確な要因について高い検出性能を示した。一方，高温や強風，近傍駅波及のように影響が小さいまたは間接的な要因では性能が低下した。これは，\LLM\ が数値データから因果関係を推論する際に，影響度の大きさと閾値の明示性が重要であることを示している。

\section{考察}
需要推定の結果より，単駅の自己回帰による予測（Last-hour, Ridge）には限界があるが，\GCN\ による近傍特徴を加えることで誤差が大幅に減少した。これは，鉄道需要の短期変動が単駅の自己履歴だけでなく，周辺駅の同時変動と強く結び付いていることを示している。したがって，本研究で構成した機能的フローネットワークは，説明生成以前の予測段階においても有効な構造表現となっている。

また，ネットワーク診断結果では，エッジ上の平均類似度およびラグ付き相互相関が高く，最大相関がほぼ $\pm1$ band に収まった。これは，機能的フローネットワークが駅間の共変動を妥当に表現しており，\GCN\ の近傍集約仮定と整合していることを意味する。単なる地理隣接ではなく，利用パターンの連動を表す構造を採用したことが，予測と説明の両面で効いていると考えられる。

説明生成では，外生要因と近傍駅変動情報を入力することで，降雨（F1: 0.718）や週末（F1: 0.696）といった明瞭な変動要因について高い検出性能を示した。これは，\LLM\ が内部知識だけに頼るのではなく，観測事実に基づく候補要因を参照することで，説明の妥当性を高められることを意味する。

加えて，駅間波及情報を追加した条件では近傍駅変動に関する特徴量を与えたものの，\LLM\ 単体では neighbor\_spillover ラベルの検出には至らなかった（NeighborAcc: 0.0\%）。これは，数値的な近傍変動のみでは \LLM\ が因果関係として認識することが難しく，\GCN\ により構造的に集約された特徴量を説明生成へ接続する設計の重要性を示唆している。

既存研究との比較では，交通流予測に \LLM\ による言語説明を付与する研究~\cite{guo2024towards} が存在するものの，観測した外部要因を明示的な根拠候補として入力する構成とはなっていない。本研究の枠組みは，許可ラベル集合と根拠スニペットを \LLM\ に与えることで，一般知識から派生した根拠不明の補完を抑制し，観測事実に基づく説明の妥当性を定量評価できる点で差別化される。

本研究の重要な点は，予測精度の向上と説明の成立性を同一パイプラインで扱ったことである。数値的な構造を無視して生成された説明ではなく，予測段階で得られた近傍情報と外部要因を説明生成へ直接接続したことで，説明が観測事実から大きく逸脱しにくい構成となった。特に，数値予測の改善そのものが説明対象の安定化に寄与している点は，予測と説明を別問題として扱わない本研究の意義である。

一方，本研究の評価は模擬環境に基づくため，一般化可能性には限界がある。実運用では，運行障害，イベント，振替輸送，SNS上の話題など，多様な要因が複雑に重なり合う。そのため今後は，実データを用いた検証，要因種類の拡張，および説明文と根拠の整合性を自動的に検証する仕組みの導入が必要である。また，近接ヒントをどの程度強く与えるかによって説明が変化する可能性もあるため，実データ条件では適合率と再現率のバランスも再検討する必要がある。

\section{むすび}
本稿では，路線図と気象情報を統合し，\GCN\ による需要推定と \LLM\ Agent による根拠付き説明生成を統一的に扱う枠組みを提案した。機能的フローネットワークに基づく近傍特徴を取り込むことで，需要推定誤差を0.5331から0.3673へ改善した。また，気象要因および駅間波及情報を説明生成へ与えることで，降雨・週末の要因検出で高いF1（0.718, 0.696）を達成した。提案手法は，予測精度向上と説明の妥当性を同時に実現する新枠組みとして，鉄道運用支援分野への応用が期待される。今後は，実データによる検証，外部要因の拡張，説明整合性の自動検証を進める。

\ack
本研究は，JSPS科研費（24K05080）および令和７年度大分大学研究力強化推進プロジェクト（BURST），令和７年度大分大学理工学部研究支援経費の助成を受けたものです。
\bibliography{reference}
\bibliographystyle{sieicej}

\end{document}

%\bibitem{kamarianakis2005space}
%Y.~Kamarianakis and P.~Prastacos,
%``Space-time modeling of traffic flow,''
%\textit{Computers \& Geosciences}, vol.~31, no.~2, pp.~119--133, 2005.

%\bibitem{7508408}
%S.~Hochreiter and J.~Schmidhuber,
%``Long short-term memory,''
%\textit{Neural Computation}, vol.~9, no.~8, pp.~1735--1780, 1997.
%
%\bibitem{sherstinsky2020fundamentals}
%A.~Sherstinsky,
%``Fundamentals of recurrent neural network (RNN) and long short-term memory (LSTM) network,''
%\textit{Physica D}, vol.~404, 132306, 2020.
%
%\bibitem{yu2019review}
%Y.~Yu, X.~Si, C.~Hu, and J.~Zhang,
%``A review of recurrent neural networks: LSTM cells and network architectures,''
%\textit{Neural Computation}, vol.~31, no.~7, pp.~1235--1270, 2019.
%
%\bibitem{vaswani2017attention}
%A.~Vaswani, N.~Shazeer, N.~Parmar, J.~Uszkoreit, L.~Jones, A.~N. Gomez,
%{\L}.~Kaiser, and I.~Polosukhin,
%``Attention is all you need,''
%in \textit{Proc. NeurIPS}, pp.~5998--6008, 2017.
%
%\bibitem{kipf2016semi}
%T.~N. Kipf and M.~Welling,
%``Semi-supervised classification with graph convolutional networks,''
%\textit{arXiv preprint arXiv:1609.02907}, 2016.
%
%\bibitem{defferrard2016convolutional}
%M.~Defferrard, X.~Bresson, and P.~Vandergheynst,
%``Convolutional neural networks on graphs with fast localized spectral filtering,''
%in \textit{Proc. NeurIPS}, pp.~3844--3852, 2016.
%
%\bibitem{li2017diffusion}
%Y.~Li, R.~Yu, C.~Shahabi, and Y.~Liu,
%``Diffusion convolutional recurrent neural network: Data-driven traffic forecasting,''
%\textit{arXiv preprint arXiv:1707.01926}, 2017.
%
%\bibitem{10.5555/3304222.3304273}
%B.~Yu, H.~Yin, and Z.~Zhu,
%``Spatio-temporal graph convolutional networks: A deep learning framework for traffic forecasting,''
%in \textit{Proc. IJCAI}, pp.~3634--3640, 2018.
%
%\bibitem{9363197}
%J.~Guo, X.~Liu, and Y.~Chen,
%``AST-GCN: Attribute-augmented spatiotemporal graph convolutional network for traffic forecasting,''
%\textit{IEEE Access}, vol.~9, pp.~40327--40337, 2021.
%
%\bibitem{wu2019graph}
%Z.~Wu, S.~Pan, G.~Long, J.~Jiang, X.~Chang, and C.~Zhang,
%``Graph WaveNet for deep spatial-temporal graph modeling,''
%in \textit{Proc. IJCAI}, pp.~1907--1913, 2019.
%
%\bibitem{bai2020adaptive}
%L.~Bai, L.~Yao, C.~Li, X.~Wang, and C.~Wang,
%``Adaptive graph convolutional recurrent network for traffic forecasting,''
%in \textit{Proc. NeurIPS}, vol.~33, pp.~17804--17815, 2020.
%
%\bibitem{zhou2017impacts}
%M.~Zhou, X.~Wang, M.~Li, et al.,
%``Impacts of weather on public transport ridership: Results from mining data from different sources,''
%\textit{Transportation Research Part C}, vol.~75, pp.~17--29, 2017.
%
%\bibitem{lv2015traffic}
%Y.~Lv, Y.~Duan, W.~Kang, Z.~Li, and F.~Y. Wang,
%``Traffic flow prediction with big data: A deep learning approach,''
%\textit{IEEE Transactions on Intelligent Transportation Systems}, vol.~16, no.~2, pp.~865--873, 2015.
%
%\bibitem{zhang2019deep}
%J.~Zhang, Y.~Zheng, and D.~Zhu,
%``A deep learning system for spatio-temporal traffic forecasting,''
%\textit{IEEE Transactions on Intelligent Transportation Systems}, vol.~20, no.~10, pp.~3707--3717, 2019.
%
%\bibitem{wang2018metapath}
%Y.~Wang, Y.~Zheng, and X.~Xie,
%``Metapath-based heterogeneous spatio-temporal data mining for urban traffic prediction,''
%in \textit{Proc. ACM SIGKDD}, pp.~1245--1253, 2018.
%
%\bibitem{velickovic2017graph}
%P.~Veli{\v{c}}kovi{\'c}, G.~Cucurull, A.~Casanova, A.~Romero, P.~Li{\`o}, and Y.~Bengio,
%``Graph attention networks,''
%\textit{arXiv preprint arXiv:1710.10903}, 2017.
%  bbbhh
%\bibitem{guo2024towards}
%S.~Guo, Y.~Chen, H.~Xu, et al.,
%``Towards explainable traffic flow prediction with large language models,''
%\textit{Communications in Transportation Research}, vol.~4, 100123, 2024.
%
%\bibitem{monje2022deep}
%P.~Monje, J.~A. Mena, and A.~G. Rivas,
%``Deep learning and explainable artificial intelligence for bus passenger demand prediction,''
%\textit{Applied Sciences}, vol.~12, no.~18, 2022.
%
%\bibitem{lin2004rouge}
%C.~Y. Lin,
%``ROUGE: A package for automatic evaluation of summaries,''
%in \textit{Proc. ACL Workshop}, pp.~74--81, 2004.
%
%\bibitem{scorato2023bertscore}
%E.~Scorato, M.~Valero, D.~Alfonso, et al.,
%``BERTScore: Evaluating text generation with BERT,''
%\textit{Transactions of the Association for Computational Linguistics}, vol.~11, pp.~1143--1158, 2023.
%
%\bibitem{yuan2022explainable}
%J.~Yuan, Y.~Li, H.~Zhou, et al.,
%``Explainable freight traffic forecasting with attention-based neural networks,''
%\textit{IEEE Transactions on Neural Networks and Learning Systems}, vol.~33, no.~12, pp.~4358--4370, 2022.
%
%\bibitem{zheng2015deep}
%Y.~Zheng,
%\textit{Deep Learning for Passenger Flow Prediction: Methods, Applications, and Challenges}.
%Springer, 2015.
